\documentclass[12pt,a4paper]{report}

% ── Packages ───────────────────────────────────────────────
\usepackage[a4paper, margin=2.5cm]{geometry}
\usepackage{listings}
\usepackage{xcolor}
\usepackage{setspace}
\usepackage{hyperref}
\usepackage{parskip}
\usepackage{caption}
\usepackage{float}

\onehalfspacing

\hypersetup{colorlinks=true, linkcolor=black, urlcolor=blue}

% ── Code listing style ─────────────────────────────────────
\definecolor{codebg}{RGB}{245,245,245}
\definecolor{codeframe}{RGB}{200,200,200}
\definecolor{codecomment}{RGB}{80,120,80}
\definecolor{codekeyword}{RGB}{0,70,160}
\definecolor{codestring}{RGB}{160,30,30}
\definecolor{codenumber}{RGB}{120,120,120}

\lstdefinestyle{agristyle}{
  backgroundcolor=\color{codebg},
  frame=single,
  framerule=0.5pt,
  rulecolor=\color{codeframe},
  basicstyle=\ttfamily\footnotesize,
  keywordstyle=\color{codekeyword}\bfseries,
  stringstyle=\color{codestring},
  commentstyle=\color{codecomment}\itshape,
  numberstyle=\tiny\color{codenumber},
  numbers=left,
  numbersep=8pt,
  stepnumber=1,
  showstringspaces=false,
  breaklines=true,
  breakatwhitespace=false,
  tabsize=4,
  captionpos=b,
  aboveskip=12pt,
  belowskip=8pt,
  xleftmargin=18pt,
}

% Python style
\lstdefinestyle{python}{
  style=agristyle,
  language=Python,
  morekeywords={torch,nn,F,transforms,models,optim,Path,json,os},
}

% TypeScript style (treated as JavaScript)
\lstdefinestyle{typescript}{
  style=agristyle,
  language=JavaScript,
  morekeywords={const,let,type,interface,export,import,from,async,await,
                Promise,Router,Request,Response,NextFunction,void,string,
                boolean,number,null,undefined,Record,Array,Set,true,false},
}

% SQL style
\lstdefinestyle{sql}{
  style=agristyle,
  language=SQL,
  morekeywords={INSERT,INTO,VALUES,ON,CONFLICT,DO,UPDATE,SET,WHERE,
                CREATE,TABLE,IF,NOT,EXISTS,REFERENCES,DEFAULT,UUID,TEXT,
                BOOLEAN,NUMERIC,TIMESTAMPTZ,NOW},
}

% ──────────────────────────────────────────────────────────
\begin{document}
% ──────────────────────────────────────────────────────────

\chapter*{Appendix A --- Source Code (Key Modules)}
\addcontentsline{toc}{chapter}{Appendix A --- Source Code}

This appendix presents the core source code of the AgriSmart platform across six key modules spanning the ML pipeline, backend API, and data layer. Non-essential boilerplate and import statements are shown where relevant but truncated comments are indicated with \texttt{...} where space requires.

\tableofcontents

\newpage

% ============================================================
\section*{A.1 \quad ML Inference Server --- \texttt{ml/serve.py}}
\addcontentsline{toc}{section}{A.1 \quad ML Inference Server (ml/serve.py)}
% ============================================================

The Flask inference server loads the trained MobileNetV2 checkpoint from disk on startup, then serves predictions via two endpoints: \texttt{POST /predict} (accepts an image URL) and \texttt{POST /predict-upload} (accepts a multipart file). Temperature Scaling ($T = 1.8$) is applied to logits before softmax to prevent overconfident predictions.

\begin{lstlisting}[style=python, caption={Model loading and inference server initialisation (ml/serve.py)}]
import io, json, os, sys, requests, torch
import torch.nn.functional as F
from flask import Flask, jsonify, request
from flask_cors import CORS
from PIL import Image
from torchvision import models, transforms
import torch.nn as nn
from pathlib import Path

MODEL_PATH     = Path(__file__).parent / 'outputs' / 'best_model.pth'
CLASS_MAP_PATH = Path(__file__).parent / 'outputs' / 'class_mapping.json'
PORT           = int(os.environ.get('INFER_PORT', 5001))
IMG_SIZE       = 224
DEVICE = (torch.device('cuda')  if torch.cuda.is_available()  else
          torch.device('mps')   if torch.backends.mps.is_available() else
          torch.device('cpu'))

with open(CLASS_MAP_PATH) as f:
    CLASS_MAPPING: dict[str, str] = json.load(f)   # {"0": "Potato___Early_blight", ...}
NUM_CLASSES = len(CLASS_MAPPING)

def build_model(arch: str, num_classes: int):
    if arch == 'mobilenetv2':
        m = models.mobilenet_v2(weights=None)
        in_features = m.classifier[1].in_features
        m.classifier = nn.Sequential(
            nn.Dropout(0.3),
            nn.Linear(in_features, 512),
            nn.ReLU(),
            nn.Dropout(0.2),
            nn.Linear(512, num_classes),
        )
    return m

# Load checkpoint once at startup (prevents per-request latency)
checkpoint   = torch.load(MODEL_PATH, map_location=DEVICE, weights_only=False)
arch         = checkpoint.get('arch', 'mobilenetv2')
saved_classes = checkpoint.get('class_names', list(CLASS_MAPPING.values()))
model = build_model(arch, len(saved_classes))
model.load_state_dict(checkpoint['model_state_dict'])
model.to(DEVICE).eval()
\end{lstlisting}

\begin{lstlisting}[style=python, caption={Temperature-scaled prediction function and Flask endpoints (ml/serve.py)}]
preprocess = transforms.Compose([
    transforms.Resize(int(IMG_SIZE * 1.14)),
    transforms.CenterCrop(IMG_SIZE),
    transforms.ToTensor(),
    transforms.Normalize(mean=[0.485, 0.456, 0.406],
                         std=[0.229, 0.224, 0.225]),
])

def predict_from_image(img: Image.Image) -> dict:
    """Run inference with Temperature Scaling (T=1.8) to calibrate
    overconfident predictions into a realistic 75-95% confidence range."""
    TEMPERATURE = 1.8

    tensor = preprocess(img.convert('RGB')).unsqueeze(0).to(DEVICE)
    with torch.no_grad():
        logits       = model(tensor)
        scaled_logits = logits / TEMPERATURE     # Temperature scaling
        probs         = F.softmax(scaled_logits, dim=1).squeeze(0).cpu().tolist()

    predictions = [
        {'label': saved_classes[i], 'confidence': round(probs[i], 6)}
        for i in range(len(saved_classes))
    ]
    predictions.sort(key=lambda x: x['confidence'], reverse=True)
    top = predictions[0]
    return {'top_label': top['label'],
            'top_confidence': top['confidence'],
            'predictions': predictions}

app = Flask(__name__)
CORS(app)

@app.get('/health')
def health():
    return jsonify({'status': 'ok', 'classes': len(saved_classes), 'arch': arch})

@app.post('/predict')
def predict():
    body      = request.get_json(force=True, silent=True) or {}
    image_url = body.get('imageUrl', '').strip()
    if not image_url:
        return jsonify({'error': 'imageUrl is required'}), 400
    resp = requests.get(image_url, timeout=15)
    img  = Image.open(io.BytesIO(resp.content))
    return jsonify({'success': True, **predict_from_image(img)})

@app.post('/predict-upload')
def predict_upload():
    """Accepts multipart file upload (used by Express API server)."""
    file   = request.files['file']
    img    = Image.open(file.stream)
    return jsonify({'success': True, **predict_from_image(img)})

if __name__ == '__main__':
    app.run(host='0.0.0.0', port=PORT, debug=False)
\end{lstlisting}

\newpage

% ============================================================
\section*{A.2 \quad Model Training Pipeline --- \texttt{ml/train.py}}
\addcontentsline{toc}{section}{A.2 \quad Model Training Pipeline (ml/train.py)}
% ============================================================

The training pipeline implements two-phase transfer learning: Phase 1 trains only the custom classifier head with a frozen backbone; Phase 2 unfreezes the last feature blocks for fine-tuning.  Class-weighted loss and weighted random sampling counter severe class imbalance in the PlantVillage dataset.

\begin{lstlisting}[style=python, caption={Data augmentation pipeline for real-world generalization (ml/train.py)}]
def get_transforms(img_size: int):
    """Aggressive augmentation strategy — forces the model to learn disease
    texture features rather than background colour cues.

    1. RandomResizedCrop (scale 0.4-1.0): partial leaves, locality learning
    2. RandomPerspective: varied camera angles in the field
    3. Heavy ColorJitter (hue=0.2): breaks background colour cues
    4. RandomGrayscale: removes colour, forces texture-only learning
    5. GaussianBlur: robustness to camera-phone shake/blur
    6. RandomErasing: occludes random background patches post-ToTensor
    """
    train_transform = transforms.Compose([
        transforms.RandomResizedCrop(img_size, scale=(0.4, 1.0), ratio=(0.75, 1.33)),
        transforms.RandomHorizontalFlip(),
        transforms.RandomVerticalFlip(),
        transforms.RandomRotation(30),
        transforms.RandomPerspective(distortion_scale=0.3, p=0.4),
        transforms.ColorJitter(brightness=0.4, contrast=0.4,
                               saturation=0.4, hue=0.2),
        transforms.RandomGrayscale(p=0.05),
        transforms.ToTensor(),
        transforms.Normalize(mean=[0.485, 0.456, 0.406],
                             std=[0.229, 0.224, 0.225]),
        transforms.RandomErasing(p=0.3, scale=(0.02, 0.15),
                                 ratio=(0.3, 3.3), value='random'),
        transforms.GaussianBlur(kernel_size=3, sigma=(0.1, 1.5)),
    ])
    val_transform = transforms.Compose([
        transforms.Resize(int(img_size * 1.14)),
        transforms.CenterCrop(img_size),
        transforms.ToTensor(),
        transforms.Normalize(mean=[0.485, 0.456, 0.406],
                             std=[0.229, 0.224, 0.225]),
    ])
    return train_transform, val_transform
\end{lstlisting}

\begin{lstlisting}[style=python, caption={Two-phase training loop with early stopping (ml/train.py)}]
def build_model(arch: str, num_classes: int):
    """Phase 1: backbone frozen. Phase 2: partial unfreeze via unfreeze_backbone()."""
    if arch == 'mobilenetv2':
        model = models.mobilenet_v2(weights=models.MobileNet_V2_Weights.IMAGENET1K_V1)
        for param in model.features.parameters():
            param.requires_grad = False            # freeze backbone
        in_features = model.classifier[1].in_features
        model.classifier = nn.Sequential(
            nn.Dropout(0.3), nn.Linear(in_features, 512),
            nn.ReLU(), nn.Dropout(0.2), nn.Linear(512, num_classes),
        )
    return model

def unfreeze_backbone(model, arch: str):
    """Phase 2: unfreeze last 4 MobileNetV2 feature blocks (14-18)."""
    if arch == 'mobilenetv2':
        for i, block in enumerate(model.features):
            if i >= 14:
                for param in block.parameters():
                    param.requires_grad = True
    return model

# --- Training loop (simplified) ---
for epoch in range(1, args.epochs + 1):
    if epoch == args.phase1_epochs + 1:          # Phase 2 transition
        model = unfreeze_backbone(model, args.model)
        optimizer = optim.Adam(
            filter(lambda p: p.requires_grad, model.parameters()),
            lr=args.lr_finetune, weight_decay=1e-4)

    train_loss, train_acc = train_one_epoch(model, train_loader,
                                            criterion, optimizer, device)
    val_loss, val_acc, _, _ = evaluate(model, val_loader, criterion, device)
    scheduler.step(val_loss)

    if val_acc > best_val_acc:                   # Save best checkpoint
        best_val_acc = val_acc
        torch.save({
            'epoch': epoch,
            'model_state_dict': model.state_dict(),
            'val_acc': val_acc,
            'class_names': class_names,
            'arch': args.model,
        }, os.path.join(args.output_dir, 'best_model.pth'))

    if epochs_without_improvement >= args.patience:
        break                                    # Early stopping
\end{lstlisting}

\newpage

% ============================================================
\section*{A.3 \quad Multi-Model Inference Router --- \texttt{api/routes/roboflow.ts}}
\addcontentsline{toc}{section}{A.3 \quad Multi-Model Inference Router (api/routes/roboflow.ts)}
% ============================================================

This Express router is the central inference orchestrator. It receives a \texttt{POST /api/roboflow/infer} request from the React frontend, routes it to the correct AI model (local Flask or one of two HuggingFace Spaces), normalises the response format, and then calls the \texttt{buildDecision()} engine with the top prediction and a matching \texttt{spray\_recipes} database record.

\begin{lstlisting}[style=typescript, caption={Model registry and HuggingFace response normaliser (api/routes/roboflow.ts)}]
const LOCAL_INFER_URL = process.env.LOCAL_INFER_URL ?? 'http://localhost:5001/predict-upload'
const HF_INSECT_URL   = 'https://SanchaiKB-Insect-Classification-Model.hf.space/predict'
const HF_RICE_URL     = 'https://sanchaikb-fertilizer-model.hf.space/predict'

const ALLOWED_MODELS = new Set([
  'agrismart/1',   // Local MobileNetV2 -- Tomato + Potato (7 classes)
  'hf-insect/1',   // HuggingFace EfficientNet-B0 Insect Classifier (5 classes)
  'hf-rice/1',     // HuggingFace EfficientNet-B0 Rice Disease Model (6 classes)
])

/** Calls a HuggingFace Space and normalises the 3 possible response formats
 *  into a unified { predictions: [{label, confidence}] } structure. */
async function callHFEndpoint(hfUrl, imageUrl) {
  const { blob, filename } = await fetchImageAsBlob(imageUrl)
  const form = new FormData()
  form.append('file', blob, filename)

  const r = await fetch(hfUrl, { method: 'POST', body: form,
                                  signal: AbortSignal.timeout(45_000) })
  const data = await r.json()

  // --- Format 1: Rice model returns { scores: [{label, confidence}] }
  if (Array.isArray(data.scores)) {
    const preds = data.scores.map(s => ({ label: s.label, confidence: s.confidence }))
    preds.sort((a, b) => b.confidence - a.confidence)
    return { predictions: preds, top_label: preds[0]?.label, ...data }
  }

  // --- Format 2: EfficientNet insect model returns { all_scores: {label: num} }
  if (data.all_scores && typeof data.all_scores === 'object') {
    const preds = Object.entries(data.all_scores)
      .map(([label, confidence]) => ({ label, confidence }))
      .sort((a, b) => b.confidence - a.confidence)
    return { predictions: preds, top_label: preds[0]?.label, ...data }
  }

  // --- Format 3: Legacy top-1 only { prediction, confidence }
  return {
    predictions: [{ label: String(data.prediction), confidence: Number(data.confidence) }],
    top_label: String(data.prediction), ...data,
  }
}
\end{lstlisting}

\begin{lstlisting}[style=typescript, caption={Main inference route with model routing and decision engine (api/routes/roboflow.ts)}]
router.post('/infer', requireAuth, async (req, res) => {
  const { modelId, imageUrl, plantId } = req.body

  if (!ALLOWED_MODELS.has(modelId))
    return res.status(400).json({ error: `Unsupported modelId: ${modelId}` })
  if (!imageUrl || !/^https?:\/\//.test(imageUrl))
    return res.status(400).json({ error: 'Invalid imageUrl' })

  let raw
  // ── Route to correct model ──────────────────────────────────────────
  if (modelId === 'agrismart/1') {
    const { blob, filename } = await fetchImageAsBlob(imageUrl)
    const form = new FormData()
    form.append('file', blob, filename)
    const r    = await fetch(LOCAL_INFER_URL, { method: 'POST', body: form })
    const data = await r.json()
    raw = { predictions: data.predictions, top_label: data.top_label }

  } else if (modelId === 'hf-insect/1') {
    raw = await callHFEndpoint(HF_INSECT_URL, imageUrl)

  } else {   // hf-rice/1
    raw = await callHFEndpoint(HF_RICE_URL, imageUrl)
  }

  // ── Spray recipe lookup from Supabase ───────────────────────────────
  const sorted    = [...raw.predictions].sort((a, b) => b.confidence - a.confidence)
  const top       = sorted[0] ?? { label: 'unknown', confidence: 0 }
  const supabase  = createSupabaseClient(req.accessToken)

  const { data: recipe } = await supabase
    .from('spray_recipes')
    .select('*')
    .eq('model_id', modelId)
    .eq('class_label', top.label)
    .limit(1)
    .maybeSingle()

  // ── Build decision ──────────────────────────────────────────────────
  const decision = buildDecision({ modelId, raw: { predictions: raw.predictions }, recipe })
  res.status(200).json({ success: true, raw, decision, plantId })
})
\end{lstlisting}

\newpage

% ============================================================
\section*{A.4 \quad Decision Engine --- \texttt{api/lib/decisionEngine.ts}}
\addcontentsline{toc}{section}{A.4 \quad Decision Engine (api/lib/decisionEngine.ts)}
% ============================================================

The decision engine implements a deterministic three-state rules logic. It receives the raw model output and the matching \texttt{spray\_recipes} database row, and returns a structured \texttt{Decision} object that drives the UI badge, recommendation card, and spray schedule creation.

\begin{lstlisting}[style=typescript, caption={Three-state decision logic — Healthy / Action Required / Review (api/lib/decisionEngine.ts)}]
export function buildDecision(args: {
  modelId: string
  raw: RoboflowRawResponse
  recipe: SprayRecipe | null
}): Decision {
  const preds   = normalizePredictions(args.raw)
  const { top, second } = getTopTwo(preds)
  const margin  = second ? top.confidence - second.confidence : 1

  // Low confidence or ambiguous (margin too small between top-2)
  const isUncertain = top.confidence < 0.6 || margin < 0.08

  // ── State 1: HEALTHY (no action) ────────────────────────────────────
  // Matches: 'Healthy', 'Healthy Crop', 'Tomato___Healthy', '_healthy', etc.
  const isHealthy =
    top.label.toLowerCase() === 'healthy' ||
    top.label.toLowerCase() === 'healthy crop' ||
    top.label.toLowerCase().endsWith('___healthy') ||
    top.label.toLowerCase().endsWith('_healthy')

  if (isHealthy) {
    return {
      spray: false, actionType: 'none', label: top.label,
      confidence: top.confidence,
      recommendation: 'Healthy',
      dosage: '-',
      notes: 'Crop appears healthy. No treatment required.',
      reason: 'Healthy class detected.', ruleVersion: RULE_VERSION,
    }
  }

  // ── State 2: MISSING RECIPE ──────────────────────────────────────────
  if (!args.recipe || !args.recipe.enabled) {
    return {
      spray: false, actionType: 'none', label: top.label,
      confidence: top.confidence,
      recommendation: 'No product configured',
      dosage: '-',
      notes: 'No recipe configured for this class yet.',
      reason: 'Missing spray recipe mapping.', ruleVersion: RULE_VERSION,
    }
  }

  // ── State 3: LOW CONFIDENCE (review) ────────────────────────────────
  if (isUncertain || top.confidence < args.recipe.min_confidence) {
    return {
      spray: false, actionType: 'none', label: top.label,
      confidence: top.confidence,
      recommendation: 'Low confidence',
      dosage: '-',
      notes: 'Retake a clearer image or try different lighting.',
      reason: `Confidence ${top.confidence.toFixed(2)} below threshold.`,
      ruleVersion: RULE_VERSION,
    }
  }

  // ── State 4: ACTION REQUIRED ─────────────────────────────────────────
  const actionType = args.recipe.action_type === 'fertilizer' ? 'fertilizer'
                   : args.recipe.action_type === 'pesticide'  ? 'pesticide'
                   : 'none'

  return {
    spray: actionType !== 'none',
    actionType, label: top.label, confidence: top.confidence,
    recommendation: args.recipe.recommendation,
    dosage:  args.recipe.dosage || 'Follow label instructions.',
    notes:   args.recipe.notes  || '',
    reason:  `Matched recipe for ${top.label}.`,
    ruleVersion: RULE_VERSION,
  }
}
\end{lstlisting}

\newpage

% ============================================================
\section*{A.5 \quad JWT Authentication Middleware --- \texttt{api/middleware/auth.ts}}
\addcontentsline{toc}{section}{A.5 \quad JWT Authentication Middleware (api/middleware/auth.ts)}
% ============================================================

Every Express route in the AgriSmart API is protected by the \texttt{requireAuth} middleware. It extracts the \texttt{Bearer} token from the \texttt{Authorization} header and validates it against Supabase Auth. If valid, the user ID and token are attached to the request object for downstream use by route handlers.

\begin{lstlisting}[style=typescript, caption={JWT validation middleware for all protected API routes (api/middleware/auth.ts)}]
import type { Request, Response, NextFunction } from 'express'
import { createSupabaseClient } from '../lib/supabaseAdmin.js'

export interface AuthenticatedRequest extends Request {
  userId?:      string   // Supabase user UUID -- set on successful auth
  accessToken?: string   // Raw JWT -- used to create user-scoped Supabase client
}

export function requireAuth(
  req: AuthenticatedRequest,
  res: Response,
  next: NextFunction
): void {
  const authHeader = req.headers.authorization
  if (!authHeader || !authHeader.startsWith('Bearer ')) {
    res.status(401).json({ success: false,
                           error: 'Missing or invalid Authorization header' })
    return
  }

  const token    = authHeader.slice(7)          // Strip "Bearer " prefix
  const supabase = createSupabaseClient(token)

  supabase.auth.getUser(token).then(({ data, error }) => {
    if (error || !data.user) {
      res.status(401).json({ success: false, error: 'Invalid or expired token' })
      return
    }
    // Attach user context to request for downstream route handlers
    req.userId      = data.user.id
    req.accessToken = token
    next()
  }).catch(() => {
    res.status(401).json({ success: false, error: 'Authentication failed' })
  })
}
\end{lstlisting}

\newpage

% ============================================================
\section*{A.6 \quad Spray Schedule API --- \texttt{api/routes/schedules.ts}}
\addcontentsline{toc}{section}{A.6 \quad Spray Schedule API (api/routes/schedules.ts)}
% ============================================================

The spray schedules route manages the lifecycle of treatment plans created after each diagnosis. The key endpoint is \texttt{PATCH /:id/complete}, which increments the completed application counter and recalculates the next spray date, automatically marking the schedule as \texttt{completed} once all applications are done.

\begin{lstlisting}[style=typescript, caption={Spray schedule creation and mark-complete endpoints (api/routes/schedules.ts)}]
// POST / -- Create a new schedule from a diagnosis result
router.post('/', requireAuth, async (req, res) => {
  const { diagnosisId, plantId, productName,
          dosage, intervalDays, totalApplications } = req.body

  const supabase  = createSupabaseClient(req.accessToken)
  const today     = new Date().toISOString().split('T')[0]

  const { data, error } = await supabase
    .from('spray_schedules')
    .insert({
      diagnosis_id:           diagnosisId,
      plant_id:               plantId || null,
      user_id:                req.userId,
      product_name:           productName,
      dosage:                 dosage || 'Follow label instructions',
      interval_days:          intervalDays  || 7,
      total_applications:     totalApplications || 3,
      completed_applications: 0,
      start_date:             today,
      next_spray_date:        today,
      status:                 'active',
    })
    .select('*').single()

  res.status(201).json({ success: true, schedule: data })
})

// PATCH /:id/complete -- Mark one spray application as done
router.patch('/:id/complete', requireAuth, async (req, res) => {
  const supabase = createSupabaseClient(req.accessToken)
  const { id }   = req.params

  // Fetch the current schedule to read counters
  const { data: schedule } = await supabase
    .from('spray_schedules').select('*')
    .eq('id', id).eq('user_id', req.userId).single()

  const completedApps = (schedule.completed_applications || 0) + 1
  const isComplete    = completedApps >= (schedule.total_applications || 3)

  // Recalculate next spray date from today + interval_days
  const nextDate = new Date()
  nextDate.setDate(nextDate.getDate() + (schedule.interval_days || 7))

  const { data } = await supabase
    .from('spray_schedules')
    .update({
      completed_applications: completedApps,
      next_spray_date: isComplete
        ? schedule.next_spray_date          // frozen if fully complete
        : nextDate.toISOString().split('T')[0],
      status: isComplete ? 'completed' : 'active',
    })
    .eq('id', id).select('*').single()

  res.status(200).json({ success: true, schedule: data })
})
\end{lstlisting}

\newpage

% ============================================================
\section*{A.7 \quad Master Spray Recipes --- \texttt{supabase/migrations/0008\_master\_recipes.sql}}
\addcontentsline{toc}{section}{A.7 \quad Master Spray Recipes SQL Migration}
% ============================================================

The \texttt{spray\_recipes} table is the single source of truth for all 18 treatment configurations across the three AI models. Each row maps a \texttt{(model\_id, class\_label)} pair to a chemical product, dosage, action type, confidence threshold, and special field notes.

\begin{lstlisting}[style=sql, caption={Spray recipes for Rice Disease model --- 6 classes (0008\_master\_recipes.sql, excerpt)}]
-- Rice Disease model (hf-rice/1) -- EfficientNet-B0, 6 classes
INSERT INTO spray_recipes
  (model_id, class_label, enabled, min_confidence,
   action_type, recommendation, dosage, notes)
VALUES
  ('hf-rice/1', 'Blast Disease', true, 0.60,
   'pesticide',
   'Tricyclazole (Beam) or Kasugamycin Fungicide',
   '0.6g Tricyclazole per litre of water; spray every 10-14 days',
   'Spray at tillering and booting stages. Avoid spraying during
    rain. Do not apply within 14 days of harvest.'),

  ('hf-rice/1', 'Brown Spot Disease', true, 0.60,
   'fertilizer',
   'Mancozeb or Iprodione Fungicide + Potash Fertiliser',
   '2g Mancozeb per litre; apply 60-90 kg/ha Potash',
   'Brown spot is often aggravated by potassium deficiency.
    Apply potash before fungicide for best results.'),

  ('hf-rice/1', 'Bacterial Blight Disease', true, 0.60,
   'pesticide',
   'Copper Oxychloride or Streptomycin Bactericide',
   '3g Copper Oxychloride per litre, spray every 7-10 days',
   'Remove and destroy infected leaves. Avoid overhead irrigation.
    Do not reuse water from infected fields. Practice crop rotation.'),

  ('hf-rice/1', 'False Smut Disease', true, 0.55,
   'pesticide',
   'Propiconazole or Hexaconazole Fungicide',
   '1ml Propiconazole per litre, apply at heading stage',
   'Apply preventively at boot leaf stage. Avoid excessive nitrogen.
    Infected grains are toxic -- do not use as seed for next season.'),

  ('hf-rice/1', 'NeckBlast', true, 0.65,
   'pesticide',
   'Tricyclazole (Beam) or Kasugamycin Fungicide',
   '0.6g per litre of water, spray at neck emergence',
   'Critical application window: spray at 50% heading.
    A second spray 10 days later recommended for heavy infection.'),

  ('hf-rice/1', 'Healthy Crop', true, 0.50,
   'none', 'No treatment required', '-',
   'Crop is healthy. Continue standard agronomic practices.')
ON CONFLICT (model_id, class_label)
  DO UPDATE SET
    recommendation = EXCLUDED.recommendation,
    dosage         = EXCLUDED.dosage,
    notes          = EXCLUDED.notes,
    updated_at     = NOW();
\end{lstlisting}

\begin{lstlisting}[style=sql, caption={Spray recipes for EfficientNet-B0 Insect model --- 5 pest classes (0008\_master\_recipes.sql, excerpt)}]
-- Insect/Pest model (hf-insect/1) -- EfficientNet-B0, 5 classes
INSERT INTO spray_recipes
  (model_id, class_label, enabled, min_confidence,
   action_type, recommendation, dosage, notes)
VALUES
  ('hf-insect/1', 'Rice Stem Borer', true, 0.60,
   'pesticide',
   'Chlorpyrifos or Carbofuran Granules',
   '1.5ml Chlorpyrifos per litre of water, or 1kg Carbofuran
    granules per acre at transplanting',
   'Apply granules near plant base. Avoid excessive post-application
    irrigation. Repeat in 15 days if infestation persists.'),

  ('hf-insect/1', 'Green Leafhopper', true, 0.60,
   'pesticide',
   'Imidacloprid or Thiamethoxam Insecticide',
   '0.5ml Imidacloprid per litre; spray every 7-10 days',
   'Green leafhoppers are vectors for Rice Tungro Virus.
    Early treatment is critical. Avoid spraying beneficial insects.'),

  ('hf-insect/1', 'Planthopper', true, 0.60,
   'pesticide',
   'Buprofezin or Pymetrozine Insecticide',
   '1ml Buprofezin per litre; 2 applications 10 days apart',
   'Drain field before spraying for contact action.
    Avoid consecutive use of same insecticide class (resistance risk).'),

  ('hf-insect/1', 'Rice Bug', true, 0.55,
   'pesticide',
   'Malathion or Lambda-cyhalothrin Insecticide',
   '2ml Malathion per litre; spray at milk and dough stage',
   'Critical period: heading to grain fill. Early-morning spraying
    most effective when bugs are less active.'),

  ('hf-insect/1', 'Rice Leaf Roller', true, 0.55,
   'pesticide',
   'Chlorantraniliprole (Coragen) or Flubendiamide',
   '0.3ml Coragen per litre; apply at egg hatch stage',
   'Foliar spray targeting young larvae before leaf rolling.
    Avoid spraying at flowering to protect pollinators.')
ON CONFLICT (model_id, class_label)
  DO UPDATE SET
    recommendation = EXCLUDED.recommendation,
    dosage         = EXCLUDED.dosage,
    notes          = EXCLUDED.notes,
    updated_at     = NOW();
\end{lstlisting}

% ──────────────────────────────────────────────────────────
\end{document}
% ──────────────────────────────────────────────────────────
